\documentclass[10pt,a4paper]{article}

\usepackage[T1]{fontenc}
\usepackage[utf8]{inputenc}
\usepackage{lmodern}
\usepackage[margin=2.15cm]{geometry}
\usepackage{amsmath,amssymb}
\usepackage{xcolor}
\usepackage{booktabs}
\usepackage{microtype}
\usepackage[hidelinks]{hyperref}

\definecolor{qrtblue}{HTML}{174A7E}
\definecolor{qrtlight}{HTML}{EEF4F9}
\setlength{\parindent}{0pt}
\setlength{\parskip}{0.55em}

\begin{document}

\begin{center}
  {\Large\bfseries\color{qrtblue} Econometric Formulation of the Problem}\\[0.35em]
  {\small A hierarchical panel model for predicting the sign of a future return}
\end{center}

\section*{1. Panel structure}

The data form a panel. Allocation $a=1,\ldots,A$ is observed on date
$t=1,\ldots,T$ and belongs to group $g(a)$. Let
\[
  r_{a,t+1} \quad \text{be the future return},
  \qquad
  y_{a,t}=\mathbf 1\{r_{a,t+1}>0\}
\]
be the continuous outcome and the associated binary target. At date $t$ we
observe:
\begin{itemize}
  \item $x_{a,t}$: allocation-specific information known at prediction time;
  \item $z_t$: the common market state, also known at prediction time.
\end{itemize}

The objective is to separate what is common to every observation from what is
specific to a group, an allocation, or a date.

\section*{2. General decomposition}

A convenient econometric equation is
\begin{equation}
\boxed{
\begin{aligned}
 r_{a,t+1}
 ={}& \underbrace{\mu}_{\text{global mean}}
 + \underbrace{\theta_{g(a)}}_{\text{group effect}}
 + \underbrace{\alpha_a}_{\text{allocation effect}} \\
 &+ \underbrace{z_t^{\top}\kappa}_{\text{market-state effect}}
 + \underbrace{x_{a,t}^{\top}\beta}_{\text{common slopes}}
 + \underbrace{x_{a,t}^{\top}b_a}_{\text{slope deviation}}
 + \underbrace{\varepsilon_{a,t+1}}_{\text{future shock}}.
\end{aligned}
}
\label{eq:panel}
\end{equation}

The components have the following meaning:
\begin{center}
\begin{tabular}{@{}ll@{}}
\toprule
Term & Interpretation \\
\midrule
$\mu$ & average return across the whole panel; \\
$\theta_{g(a)}$ & persistent difference between groups; \\
$\alpha_a$ & allocation $a$'s deviation from its group; \\
$z_t^\top\kappa$ & effect of the market state shared on date $t$; \\
$x_{a,t}^\top\beta$ & response to predictors common to all allocations; \\
$x_{a,t}^\top b_a$ & allocation-specific deviation from the common response; \\
$\varepsilon_{a,t+1}$ & future information not known at date $t$. \\
\bottomrule
\end{tabular}
\end{center}

Hence allocation $a$ has total intercept
\[
  \mu+\theta_{g(a)}+\alpha_a
\]
and total slope vector
\[
  \beta+b_a.
\]
This makes the decomposition easy to read: there is a common model, to which
we add progressively finer group- and allocation-level corrections.

\section*{3. The date effect and forecasting}

For an in-sample panel regression one could write an unrestricted date fixed
effect $\lambda_t$. It is useful conceptually to decompose it as
\[
  \lambda_t=z_t^\top\kappa+\nu_t,
\]
where $z_t^\top\kappa$ is explained by information observed on date $t$ and
$\nu_t$ is an unpredictable common shock. An unrestricted $\lambda_t$ cannot
be used to forecast a new date because its coefficient is not yet known.
Therefore the forecasting equation uses the observable component
$z_t^\top\kappa$; the remaining $\nu_t$ is absorbed into the forecast error.

\section*{4. From return prediction to sign prediction}

Define the systematic score, excluding the future shock, by
\[
  \eta_{a,t}
  =\mu+\theta_{g(a)}+\alpha_a+z_t^\top\kappa
   +x_{a,t}^\top\beta+x_{a,t}^\top b_a.
\]
There are two equivalent modeling routes.

\textbf{Continuous route.} Estimate Equation~\eqref{eq:panel}, predict
$\widehat r_{a,t+1}=\widehat\eta_{a,t}$, and use
\[
  \widehat y_{a,t}=\mathbf 1\{\widehat\eta_{a,t}>0\}.
\]
If $\varepsilon_{a,t+1}\sim\mathcal N(0,\sigma^2)$, then
\[
  \Pr(y_{a,t}=1\mid x_{a,t},z_t,a)
  =\Phi\!\left(\frac{\eta_{a,t}}{\sigma}\right).
\]

\textbf{Binary route.} Specify the sign probability directly as
\begin{equation}
  \Pr(y_{a,t}=1\mid x_{a,t},z_t,a)
  =F(\eta_{a,t}),
\label{eq:binary}
\end{equation}
where $F$ is the logistic cdf (logit) or the standard normal cdf (probit).
With either symmetric link, the default $50\%$ decision rule is equivalent to
$\widehat\eta_{a,t}>0$.

\section*{5. Partial pooling and identification}

Estimating an unrelated coefficient for every allocation is noisy. A
hierarchical specification shares information across the panel:
\[
  \theta_g\sim\mathcal N(0,\sigma_g^2),\qquad
  \alpha_a\sim\mathcal N(0,\sigma_a^2),\qquad
  b_a\sim\mathcal N(0,\Sigma_b).
\]
Small or weakly observed groups and allocations are then shrunk toward the
common model, while well-observed ones may differ when the data support it.
The same principle can be implemented with penalized fixed effects.

For a pure fixed-effect formulation, impose normalizations such as
\[
  \sum_g\theta_g=0,
  \qquad
  \sum_{a:g(a)=g}\alpha_a=0,
  \qquad
  \sum_a b_a=0,
\]
so that the global, group, and allocation components are separately identified.

\vfill
\begin{center}
\colorbox{qrtlight}{\parbox{0.91\linewidth}{\small
\textbf{Core model.} Future returns are decomposed into a global component,
nested group and allocation effects, an observable market-state effect, common
predictor slopes, allocation-specific slope deviations, and an unpredictable
shock. The classification target is simply the sign of this latent return.}}
\end{center}

\end{document}
